% =========================================================
\chapter{C01 --- Les signaux et les sources}
% =========================================================

\begin{rembox}[title=En bref]
\textbf{De quoi parle ce chapitre ?} Des \textbf{signaux électriques} (tension $U$ et courant
$I$, leurs formes), des \textbf{sources} qui les créent (idéales vs réelles, de tension et de
courant), des \textbf{conventions} (fléchage, signe de la puissance) et de la
\textbf{puissance} : valeur moyenne et \textbf{valeur efficace}. \\
\textbf{Les deux formules du jour :} $U = R\cdot I$ (loi d'Ohm) et $P = U\cdot I$
(puissance électrique).
\end{rembox}

Les trois choses à maîtriser à la fin du chapitre (diapositive 2) :
\begin{enumerate}
  \item \textbf{Maîtriser les signaux électriques} : leur nature (tension, courant), leurs
        formes (sinus, carré, triangle, continu), les définitions clés et la convention de
        fléchage de l'EPFL, directement liée à la notion de puissance.
  \item \textbf{Sources idéales vs réelles} : comment imposer une tension ou un courant à un
        circuit ; les sources \og idéales\fg{} (impossibles en pratique) comparées aux sources
        \og réelles\fg{}, en tension et en courant.
  \item \textbf{Analyser la puissance} : comprendre la puissance efficace (ou moyenne) pour
        évaluer la consommation et la tolérance d'un circuit ; c'est de là que viennent les
        valeurs efficaces $U_{EFF}$, $I_{EFF}$ utilisées partout.
\end{enumerate}

\section{C'est quoi un signal ?}

